\documentclass[11pt]{article}
\usepackage{amsmath,amssymb,amsthm}

\begin{document}

\begin{proof}
Maximum-Norm A Posteriori Upper Bound via Green Representation and Elliptic Reconstruction for Backward Euler/FEM.

Let $\Omega\subset\mathbb{R}^d$ be bounded and let $u$ solve, for $t\in(0,T)$, the parabolic problem in the variational setting used in the statement (with operators $M$ and $K$ as in the Green-lemma framework), and let $\{u_h^j\}_{j=0}^N$ be the backward Euler/FEM solution at times $t_j=j\tau$ ($\tau_j:=t_j-t_{j-1}$). Define the time extension of the discrete solution by the piecewise linear-in-time interpolation specified in the problem summary: for $t\in I_j:=(t_{j-1},t_j)$ we set
\[
\tilde u_h(t):=(1-\theta)u_h^{j-1}+\theta u_h^j,\qquad \theta:=\frac{t-t_{j-1}}{\tau_j}.
\]
Let $u_h(t):=\tilde u_h(t)$, and define the space--time error $e(t):=u(t)-u_h(t)$. Under the regularity available for $u$ and the construction of the time extension, $e$ belongs to the Green-lemma space (at least after the standard minimal mollification/approximation if needed), so that the representation lemma applies with $w=e$. In particular, for each fixed $x\in\Omega$ the lemma yields the Green representation at the final time $T$ in the exact form
\[
e(x,T)=(\Gamma(T),e(0))+\int_0^T \langle K e(s),\Gamma(T-s)\rangle\,ds.
\]
Here $(\cdot,\cdot)$ denotes the spatial duality pairing at time $t=0$ consistent with the space where the initial datum lives, and $\langle\cdot,\cdot\rangle$ denotes the duality pairing in time--space consistent with the lemma, i.e. if $K e(s)$ is valued in the spatial dual space where the kernel $\Gamma(T-s)$ is a test function, then $\langle K e(s),\Gamma(T-s)\rangle$ is interpreted as that duality pairing (for example, $V^*$ with $V$, depending on the exact mapping of $K$ in the lemma).

Next, compute $K e$. Since $u$ satisfies the continuous equation, the parabolic defect of the reconstruction $u_h$ is exactly
\[
K e = F-(\partial_t u_h+M u_h),
\]
under the same forcing $F$ and operators as in the lemma. To express $\partial_t u_h$ we use the chosen extension: on each interval $I_j$ the function $u_h$ is affine in time, hence $\partial_t u_h$ is constant in time on $I_j$ and equals the backward difference quotient
\[
\partial_t u_h(t)=\frac{u_h^j-u_h^{j-1}}{\tau_j} \quad \text{for } t\in I_j.
\]
Therefore, for $s\in I_j$ we have the intervalwise identity
\[
K e(s)=F(s)-\Big(\frac{u_h^j-u_h^{j-1}}{\tau_j}+M\big((1-\theta)u_h^{j-1}+\theta u_h^j\big)\Big).
\]
Insert this expression into the Green representation. Since $\partial_t u_h$ and all discrete quantities are intervalwise defined, we localize the time integral by splitting at the nodes:
\[
e(x,T)=(\Gamma(T),e(0))+\sum_{j=1}^N\int_{t_{j-1}}^{t_j}\langle K e(s),\Gamma(T-s)\rangle\,ds.
\]
The standard backward Euler localization uses the time support of the affine extension together with the lemma's kernel $\Gamma$. On each $I_j$ we replace $\Gamma(T-s)$ by the appropriate interval-wise interpolation/average in time that turns the residual contribution into node-based backward Euler quantities. Concretely, choose the intervalwise test replacement required by the Green representation argument (equivalently, one uses that the representation lemma uses $\Gamma(T-s)$ exactly and then applies the same averaging used in the backward Euler consistency step); this yields that the integral over $I_j$ can be written as a sum of pairings between backward Euler time defects and kernel factors averaged over $I_j$, plus a spatial defect term that is handled by elliptic reconstruction. In particular, one arrives at the decomposition
\[
e(x,T)=(\Gamma(T),e(0)) + \sum_{j=1}^N \mathcal T_j(x) + \sum_{j=1}^N \mathcal S_j(x),
\]
where $\mathcal T_j(x)$ contains the backward Euler time residual component and $\mathcal S_j(x)$ contains the spatial residual component.

For the spatial component, define for each $j$ the elliptic reconstruction $\psi_h^j$ and defect $R^j$ exactly as in the summary. By construction, the spatial part of $K e$ on $I_j$ can be re-expressed using these quantities so that the relevant duality pairing reduces to a pairing involving $R^j-u_h^j$ (or, equivalently, $\psi_h^j$). The algebra is the standard elliptic reconstruction step: the operator $M$ acting on the affine-in-time discrete function is split into contributions controlled by the elliptic reconstruction residual, and the dependence on the affine time parameter is absorbed into intervalwise kernel averaging. Thus, for each $j$,
\[
\mathcal S_j(x)\;\text{is bounded (in duality) by a quantity controlled by the spatial reconstruction defect associated with }R^j\text{ and }u_h^j.
\]

Now bound the spatial reconstruction contribution. The given elliptic a posteriori property states that the reconstruction defect is controlled by the spatial estimator. Applying it with the identifications of the summary (the reconstructed right-hand side and the corresponding discrete solution at node $j$), we obtain an inequality of the form
\[
\sup_{x\in\Omega}|\mathcal S_j(x)|\le C\,\eta_\mathrm{ell}^j,
\]
where the norm used to define $\eta_\mathrm{ell}^j$ is compatible with the Green pairing-to-supremum conversion; specifically, the Green-function factor in $\mathcal S_j(x)$ is controlled in the same norm that appears in the elliptic estimator estimate, and the constants depend only on the coercivity/boundedness of $c/c_h$ (and elliptic constants), not on $h$ or $\tau_j$.

For the time-discretization contribution, fix $j$ and estimate $\mathcal T_j(x)$. The term $\mathcal T_j(x)$ is a duality pairing between the time defect (a backward Euler residual of the form typically $\frac{u_h^j-u_h^{j-1}}{\tau_j}-\Pi_h f^j+\cdots$, which equals the quantity described in the theorem statement) and a kernel factor involving $\Gamma(T-s)$ averaged on $I_j$. Using the available Green-function bounds, one has for the relevant kernel factor an estimate in $H^1(\Omega)$-type norms (and hence in the dual norms compatible with the time defect) of the form
\[
\|\text{(kernel factor on }I_j\text{)}\|\le C\,\omega_j,
\]
for explicit weights $\omega_j$ determined by $\Gamma$ and its time derivative. Combining this with the stability/duality inequality for the pairing gives
\[
|\mathcal T_j(x)|\le C\,\tau_j\,\|\delta_t u_h^j - f^j\|_{\ast,j},
\]
where $\delta_t u_h^j$ denotes the backward difference quotient $(u_h^j-u_h^{j-1})/\tau_j$ and $\|\cdot\|_{\ast,j}$ is precisely the computable dual norm that matches the pairing in Step 2 (the same norm in which $K e(s)$ is tested by $\Gamma$). The only use of the kernel estimates is to bound the averaged $\Gamma(T-s)$ factor on $I_j$ and its interaction with the dual norm.

The initial term is controlled similarly. At $t=0$ we have $e(0)=u_0-u_h(0)$. Since the discrete initial value is identified by the time extension used, $u_h(0)$ is known from the FEM output, hence $e(0)$ is computable in the appropriate weak sense. Using the Green-function bound for $\Gamma(T)$ (in the form supplied by the lemma, e.g. $\|\Gamma(T)\|_{1,\Omega}$ control), we obtain
\[
\sup_{x\in\Omega}|(\Gamma(T),e(0))|\le C\,\|e(0)\|_{\text{initial space}}\le C\,\|u_0-u_h(0)\|_{\text{initial space}},
\]
and the same kernel framework yields, if the final estimate is required in $\|\cdot\|_{\infty,\Omega}$, that this initial contribution fits the constant-times-computable quantity part of the estimator.

Therefore, combining the representation formula, the spatial and time bounds, and the initial bound, we have for each $x\in\Omega$,
\[
|e(x,T)|\le C\Big(\text{initial term} + \sum_{j=1}^N \tau_j\,\|\delta_t u_h^j-f^j\|_{\ast,j} + \sum_{j=1}^N \eta_\mathrm{ell}^j\Big).
\]
Since all constants in the kernel bounds and the elliptic estimator estimate are uniform in $x$, taking the supremum over $x\in\Omega$ yields
\[
\|u(T)-u_h(T)\|_{\infty,\Omega}=\sup_{x\in\Omega}|e(x,T)|\le C\Big(\text{initial term} + \sum_{j=1}^N \tau_j\,\|\delta_t u_h^j-f^j\|_{\ast,j} + \sum_{j=1}^N \eta_\mathrm{ell}^j\Big).
\]
Finally, we check consistency: each time-defect term $\delta_t u_h^j-f^j$ involves only backward Euler/FEM outputs $u_h^{j},u_h^{j-1}$ and the data $f^j$, measured in the dual norm dictated by the Green pairing; each $\eta_\mathrm{ell}^j$ depends only on $u_h^j$ and the elliptic reconstruction right-hand side as defined in the summary; the initial term depends only on $u_0$ and $u_h(0)$. Hence every quantity in the bound is computable from the backward Euler/FEM computation.

This completes the proof of the maximum-norm a posteriori upper bound for the backward Euler/FEM discretized parabolic equation, with constants depending only on the Green-decay constants and the coercivity/boundedness constants associated with $c/c_h$, as claimed in the theorem statement.
\end{proof}

\end{document}